% !TeX root = ../../Book4.tex
% 纲要标题：### 24.3 求值映射的度
% 纲要位置：几何学计划纲要.md，第 4479 行。
% 正文状态：未撰写；本文件目前是框架。

% 纲要内容开始
% 
% * 乘积复结构下的唯一球面纤维
% * 保持球内 J 的相对横截性
% 
%   正面积球面不能完全落在辛形式恰当的球内；结合简单曲线的单点注入性质，在球外扰动 J，保留球内的标准结构。
% 
% * 正则参数族和配边
% 
%   对一标记、模掉重参数化的球面模空间核验实维数为 2n；在固定类不分裂的紧性条件下，以一参数配边保持求值度。
% 
% * 求值映射的模2度为一
% * 对指定点的曲线存在
% 
%   先对正则数据由度为一取得满射，再以紧性传到需要的极限数据；不能仅以泛型点替代指定球心。
% 
% 纲要内容结束

\section{求值映射的度}
\label{g4:sec:2403}

待撰写

\subsection{乘积复结构下的唯一球面纤维}
\label{g4:subsec:2403:01}
待撰写

\subsection{保持球内 J 的相对横截性}
\label{g4:subsec:2403:02}
待撰写

\subsection{正则参数族和配边}
\label{g4:subsec:2403:03}
待撰写

\subsection{求值映射的模2度为一}
\label{g4:subsec:2403:04}
待撰写

\subsection{对指定点的曲线存在}
\label{g4:subsec:2403:05}
待撰写
% 纲要条目排为编号小节标题，并进入目录与书签；说明段仍保留于纲要注释中。
